\begingroup

\setlength{\parskip}{0.5em}
\setlength{\parindent}{0em}

\textbf{Spatial Model of Elections.}
Each voter and candidate has a position in $\mathbb{R}^2$.
The pairwise Euclidean distances are computed between
all voters and all candidates.
A voter's ballot ranks candidates in order of increasing distance.
A collection of such rankings is called a \textit{preference profile}.

\textbf{Voting Rule.}
In this context, a \textit{voting rule} is a function that takes as input a preference profile
and returns a set of one or more winning alternatives.

\textbf{Single-Winner Voting Rules.}

\begin{itemize}[topsep=0em]
\item[]
\textbf{Plurality.}
The candidate with the most first-place votes wins.

\textbf{Borda.}
Each candidate's score is
the total count of candidates ranked below them, summed across all ballots.
The candidate with the highest score wins.

\textbf{Instant Runoff Voting (IRV).}
First-place votes are tallied.
The candidate with the fewest votes is eliminated;
their supporters' votes transfer to their next favorite remaining candidate.
Repeat until only one candidate remains---they are the winner.

\end{itemize}

\textbf{Multi-Winner Voting Rules.} Let $k$ be the number of winners, i.e.,
these rules elect a committee of size $k$.

\begin{itemize}[topsep=0em]
  \item[]
  \textbf{Single Non-Transferable Vote (SNTV).}
  First-place votes are tallied.
  The $k$ highest-scoring candidates are elected.

  \textbf{Bloc Plurality.}
  Each candidate's score is the number of ballots in which
  they appear in the top $k$ positions.
  The $k$ highest-scoring candidates are elected.

  \textbf{$k$-Borda.}
  The $k$ candidates with the highest Borda scores are elected.

  \textbf{Single Transferable Vote (STV).}
  Let $n$ be the number of voters.
  First-place votes are tallied to give initial candidate scores.
  Iterate the following process:
  if any candidate's score $s$ is greater than
  $ q = \lfloor \frac{n}{k+1} \rfloor $, elect them,
  reweight each supporting voter's ballot by a factor of $\frac{s-q}{s}$,
  and transfer their support to their next-favorite remaining candidate.
  If no candidate is above the threshold, eliminate the candidate
  with the lowest score and redistribute their support.

  Only one candidate can be elected at a time:
  if multiple candidates are simultaneously above threshold,
  the candidate with the greater score is elected
  and their support is redistributed before any other candidates are elected.

\end{itemize}

\endgroup
